\chapter{Inductors and Inductive Reactance}
\label{ch:inductors}

\section{The Inductor}

An \textbf{inductor} stores energy in a magnetic field, typically by
winding a conductor into a coil. Its defining property is
\textbf{inductance} $L$, measured in \textbf{henries (H)}, defined by
how much EMF it develops in response to a changing current (a direct
application of Faraday's law, Chapter~\ref{ch:fields}):
\[
\boxed{u(t) = L\,\frac{di}{dt}}
\]
One henry is the inductance that produces one volt of induced EMF when
current changes at one ampere per second. Practical RF inductors range
from a few nanohenries (a short wire lead) to hundreds of microhenries
or more (loading coils, RF chokes).

For a simple air-core solenoid of $N$ turns, length $\ell$, and
cross-sectional area $A$ (valid when $\ell$ is much greater than the
coil's diameter),
\[
L = \frac{\mu_0 N^2 A}{\ell}.
\]
Inductance grows with the \emph{square} of the number of turns --- doubling
the turns quadruples the inductance for the same core geometry, which is
why small changes in winding count have a large effect when
hand-winding RF coils. Wrapping the coil around a high-permeability
magnetic core (ferrite or powdered iron, permeability $\mu = \mu_r
\mu_0$) multiplies the inductance by roughly $\mu_r$ for the same
winding, letting a physically small core-wound coil achieve inductance
that would otherwise require many more turns of air-core wire.

\section{Opposition to Changing Current}

Just as a capacitor opposes changes in voltage, an inductor opposes
changes in current, by Lenz's law (Chapter~\ref{ch:fields}): any attempt
to change the current through an inductor induces a back-EMF that
resists the change. Two consequences follow, mirroring the capacitor
behavior of the previous chapter:

\begin{itemize}
\item An inductor looks like a short circuit to steady DC (once current
is constant, $di/dt = 0$, so no voltage drop develops).
\item An inductor initially looks like an open circuit to a sudden
current step (opposing the abrupt change), then allows current to build
up gradually.
\end{itemize}

\begin{safetynote}[Inductive kick]
Suddenly interrupting current through an inductor (switching off a relay
coil, unplugging an energized transformer primary) forces $di/dt$ to be
very large for an instant, which can produce a large, damaging voltage
spike (``inductive kick'') and a visible spark at the switch contacts.
Flyback diodes and snubber networks are used across relay coils and
switching inductors specifically to control this spike.
\end{safetynote}

\section{Energy Stored in an Inductor}

The energy stored in an inductor carrying current $I$ is
\[
W = \frac{1}{2}LI^2.
\]

\section{Inductors in Series and Parallel}

Inductors, unlike capacitors, combine the same way resistors do (both
describe direct opposition, whether resistive or reactive, without a
reciprocal geometric relationship):

\bookfig[0.55]{inductors_series.png}{Inductors in series.}

\begin{keyidea}
\textbf{Series inductors add} (assuming no mutual coupling between them):
\[
L_{\text{total}} = L_1 + L_2 + L_3 + \cdots
\]
\textbf{Parallel inductors combine by reciprocal sum}:
\[
\frac{1}{L_{\text{total}}} = \frac{1}{L_1} + \frac{1}{L_2} + \frac{1}{L_3} + \cdots
\]
\end{keyidea}

\bookfig[0.55]{inductors_parallel.png}{Inductors in parallel.}

If two coils are wound close enough together that the magnetic field of
one links with the other (\textbf{mutual inductance}, $M$), the series
total becomes $L_{\text{total}} = L_1 + L_2 \pm 2M$, with the sign
depending on whether the windings aid or oppose each other. This
mutual-coupling effect is, in fact, the operating principle of the
transformer (Chapter~\ref{ch:transformers}).

\section{Inductive Reactance}
\label{sec:xl}

Under AC excitation, an inductor's opposition to current is
frequency-dependent, exactly as for a capacitor but with the opposite
frequency trend and the opposite phase relationship. Voltage \emph{leads}
current by $90\degree$ in an ideal inductor (the mirror image of the
capacitor case). The magnitude of this opposition is the
\textbf{inductive reactance}, $X_L$, in ohms:
\[
\boxed{X_L = 2\pi f L = \omega L}
\]

\begin{keyidea}
Inductive reactance \emph{increases} with frequency (an inductor looks
increasingly like an open circuit at high frequency) and falls to zero
as frequency approaches DC (an inductor looks like a plain wire, i.e.\ a
short circuit, at DC) --- the exact opposite frequency behavior of a
capacitor's $X_C$.
\end{keyidea}

\begin{examplebox}[Calculating reactance]
Find the reactance of a 10\,$\mu$H inductor at 14.2\,MHz.

\emph{Solution.}
\[
X_L = 2\pi f L = 2\pi \times 14.2\times10^{6} \times 10\times10^{-6} \approx 892\,\Omega.
\]
\end{examplebox}

In phasor notation, the inductor's impedance is
\[
Z_L = jX_L = j\omega L,
\]
a purely imaginary number with the opposite sign convention to the
capacitor's $Z_C = -jX_C$ --- the origin of the cancellation that occurs
at resonance, developed fully in Chapter~\ref{ch:rlc-resonance}.

\section{Practical Inductor Considerations}

\begin{itemize}
\item \textbf{Air-core coils}: no core losses, highly linear, but need
more turns (and more physical size) for a given inductance than a
core-loaded equivalent; common in high-power or high-Q RF applications
where core saturation or core loss would be problematic.
\item \textbf{Ferrite/powdered-iron core coils}: much higher inductance
per turn, widely used for RF chokes, transformers, and baluns, but
subject to \textbf{core saturation} (inductance collapses if the core's
magnetic flux density exceeds a material-specific limit) and core losses
that increase with frequency.
\item \textbf{Toroidal cores}: a closed magnetic loop (donut shape) that
confines nearly all of the magnetic field within the core, minimizing
stray coupling to nearby components --- the standard choice for compact,
well-shielded RF inductors, chokes, and transformers.
\end{itemize}

Real inductors also have parasitic series resistance (wire resistance,
increased at RF by the \textbf{skin effect}, where high-frequency current
crowds toward the conductor's surface) and parasitic parallel
capacitance between windings, which together define the coil's quality
factor $Q$ (Chapter~\ref{ch:rlc-resonance}) and its self-resonant
frequency, above which the coil behaves capacitively rather than
inductively.

\begin{quickreview}
\item $u = L\,di/dt$; $L = \mu_0 N^2 A/\ell$ for an air-core solenoid ---
inductance scales with the square of turns.
\item Inductors oppose current change: short circuit at DC, open circuit
to a sudden current step; energy stored $W = \tfrac12 LI^2$.
\item Series inductors add; parallel inductors combine by reciprocal sum.
\item $X_L = 2\pi f L$: reactance rises with frequency; voltage leads
current by $90\degree$; $Z_L = j\omega L$.
\item Magnetic cores boost inductance per turn but introduce saturation
and core-loss limits; toroids confine the field for compact, low-stray
designs.
\end{quickreview}

\begin{practiceproblems}
\item Find the reactance of a 4.7\,$\mu$H inductor at 7.1\,MHz and at
21.2\,MHz.
\item Two inductors, 15\,$\mu$H and 33\,$\mu$H, are connected in series
with negligible mutual coupling. Find the total inductance.
\item Find the energy stored in a 100\,$\mu$H inductor carrying 2\,A.
\item Explain why interrupting current through a relay coil abruptly can
produce a damaging voltage spike, and name a common protective
component used to prevent it.
\item A coil is wound on a ferrite toroid instead of air. Explain
qualitatively why fewer turns are needed to reach the same inductance,
and name one drawback of the ferrite-core approach at high power.
\end{practiceproblems}
